% ================================================================
%  PCNME: Predictive Cloud-Native Mobile Edge
%  Mathematical Methodology Paper
%  Author: Imed Eddine Benmadi
% ================================================================

\documentclass[12pt,a4paper]{article}

\usepackage[a4paper, margin=2.5cm, headheight=15pt]{geometry}
\usepackage{amsmath,amssymb,amsthm}
\usepackage{mathtools}
\usepackage{booktabs}
\usepackage{array}
\usepackage{tabularx}
\usepackage{longtable}
\usepackage{algorithm}
\usepackage{algpseudocode}
\usepackage{hyperref}
\usepackage{microtype}
\usepackage{parskip}
\usepackage{titlesec}
\usepackage{fancyhdr}
\usepackage{enumitem}
\usepackage{caption}
\usepackage{float}
\usepackage{bm}
\usepackage{xfrac}

\hypersetup{colorlinks=false, hidelinks}

\pagestyle{fancy}
\fancyhf{}
\rhead{\small PCNME: Mathematical Methodology}
\lhead{\small \thepage}
\renewcommand{\headrulewidth}{0.4pt}

\titleformat{\section}{\large\bfseries}{\thesection}{1em}{}%
  [\vspace{0.1em}\rule{\linewidth}{0.5pt}]
\titleformat{\subsection}{\normalsize\bfseries}{\thesubsection}{1em}{}
\titleformat{\subsubsection}{\normalsize\itshape\bfseries}%
  {\thesubsubsection}{1em}{}

% Theorem-like environments
\theoremstyle{definition}
\newtheorem{definition}{Definition}[section]
\newtheorem{proposition}{Proposition}[section]
\theoremstyle{remark}
\newtheorem{remark}{Remark}[section]

\newcolumntype{L}[1]{>{\raggedright\arraybackslash}p{#1}}
\newcolumntype{C}[1]{>{\centering\arraybackslash}p{#1}}

% ================================================================
\begin{document}

% ── TITLE ──────────────────────────────────────────────────────
\begin{titlepage}
\centering
\vspace*{2.5cm}
{\normalsize\textsc{Research Methodology}}\\[1.0cm]
{\LARGE\bfseries
  PCNME: Predictive Cloud-Native Mobile Edge\\[0.5em]
  A Mathematical Framework for Optimized\\[0.4em]
  Task Offloading in IoT--Fog--Cloud Systems
}\\[1.4cm]
\rule{0.55\linewidth}{0.5pt}\\[1.2cm]
{\normalsize
  \textbf{Author:} Imed Eddine Benmadi\\[0.5cm]
  \textbf{Keywords:} Task offloading, NSGA-II, MMDE,
  Execution Cost, Pareto optimisation,\\
  deep reinforcement learning, DAG,
  proactive handoff, vehicular fog computing
}\\[1.4cm]
\rule{0.55\linewidth}{0.5pt}\\[1.0cm]
{\normalsize April 2026}
\vfill
{\small\textit{eötvös loránd university - Data Science}}
\end{titlepage}

% ── ABSTRACT ────────────────────────────────────────────────────
\begin{abstract}
This paper presents the complete mathematical methodology of PCNME, the Predictive Cloud-Native Mobile Edge framework for intelligent task offloading in three-tier IoT--Fog--Cloud systems. The problem of routing computation-intensive tasks from resource-constrained IoT vehicles to fog or cloud servers is formulated as a multi-objective combinatorial optimization problem with deadline and capacity constraints. We define all system variables and derive formal expressions for \textit{\textbf{execution cost, task latency, and energy consumption}}. 

The methodology integrates five algorithmic components: an Execution Cost classifier that pre-filters tasks before optimization, an MMDE-enhanced NSGA-II multi-objective optimizer that produces Pareto-optimal offline routing plans, a behavioral cloning pipeline that seeds a Deep Q-Network agent from those plans, an online DQN agent that adapts placement decisions to live network conditions, and a GPS-driven proactive handoff predictor that prevents task failures when vehicles cross fog coverage boundaries. Every variable, equation, and parameter is defined formally and its physical interpretation is stated. The objective of this methodology section is to provide a complete, self-contained mathematical foundation from which an implementation can be derived without ambiguity.
\end{abstract}
\vspace{0.5cm}\hrule\vspace{0.7cm}
\tableofcontents
\newpage

\section{Parameters }
\label{sec:params}
% ================================================================

Table~\ref{tab:params} lists all parameters of the PCNME framework with their
symbols, default values, and justifications. Sensitivity analysis over the
key parameters $\theta$, $F$, $\mathrm{CR}$, $\gamma$, and $\omega_L$ is
recommended as part of the evaluation.

% ================================================================
% \subsection{Complete Parameter Table}
% ================================================================

Table~\ref{tab:params} collects every parameter with its symbol as used in
the mathematical formulation, its environment variable name, its default
value, and the source or physical justification for that default. The table
is organised by subsystem to match the structure of the methodology.

% --- the table spans two columns so we use longtable ---
\begingroup
\small
\renewcommand{\arraystretch}{1.3}
\begin{longtable}{@{}
    L{1.4cm}   % Symbol
    L{3.6cm}   % Description
    L{3.4cm}   % Env variable
    C{1.3cm}   % Default
    L{3.5cm}   % Source
  @{}}

\caption{Complete parameter set of the PCNME framework.}
\label{tab:params}\\

\toprule
\textbf{Symbol} &
\textbf{Description} &
\textbf{Env variable} &
\textbf{Default} &
\textbf{Source / Justification} \\
\midrule
\endfirsthead

\multicolumn{5}{l}{\small\textit{Table~\ref{tab:params} continued}}\\
\toprule
\textbf{Symbol} &
\textbf{Description} &
\textbf{Env variable} &
\textbf{Default} &
\textbf{Source / Justification} \\
\midrule
\endhead

% \midrule
% \multicolumn{5}{r}{\small\textit{Continued on next page}}\\
\endfoot

% \bottomrule
% \endlastfoot

%% ── Compute and network ─────────────────────────────────────
\multicolumn{5}{l}{\textit{Compute and network}} \\[2pt]
$\mu_k$
  & Fog node MIPS capacity
  & \texttt{FOG\_MIPS}
  & 2,000
  & Edge server benchmark \\
$\mu_c$
  & Cloud MIPS capacity
  & \texttt{CLOUD\_MIPS}
  & 8,000
  & $4{\times}$ fog (GPU server) \\
$B$
  & IoT upload bandwidth (Mbps)
  & \texttt{BANDWIDTH\_MBPS}
  & 100
  & 5G NR sub-6\,GHz \\
$B_c$
  & Fog--cloud backbone (Mbps)
  & \texttt{FOG\_CLOUD\_BW}
  & 1,000
  & Urban fibre link \\
$B_{ff}$
  & Fog--fog mesh (Mbps)
  & \texttt{FOG\_FOG\_BW}
  & 100
  & Inter-node mesh \\
$\delta_{5G}$
  & 5G one-way latency (ms)
  & \texttt{G5\_LATENCY\_MS}
  & 2
  & 3GPP Release 16 \\
$\delta_{\text{WAN}}$
  & Fog--cloud WAN delay (ms)
  & \texttt{WAN\_LATENCY\_MS}
  & 30
  & Istanbul--cloud typical \\
$\alpha$
  & WAN energy penalty factor
  & \texttt{WAN\_ENERGY\_ALPHA}
  & 1.8
  & Longer path coefficient \\
$\kappa$
  & Fog compute energy (J/MI)
  & \texttt{COMPUTE\_ENERGY\_KAPPA}
  & 0.001
  & Calibrated to ARM server \\
$P_{\text{tx}}$
  & Device transmit power (W)
  & \texttt{TX\_POWER\_WATTS}
  & 0.5
  & 5G UE typical \\

\midrule
%% ── TOF-Broker ──────────────────────────────────────────────
\multicolumn{5}{l}{\textit{TOF-Broker}} \\[2pt]
$\theta$
  & EC classification threshold (s)
  & \texttt{EC\_THRESHOLD}
  & 1.0
  & Nouri et al.~\cite{nouri2024tof} \\
$Q_{\max}$
  & Aggregator trigger depth
  & \texttt{Q\_MAX}
  & 50
  & Empirically tuned \\

\midrule
%% ── NSGA-II + MMDE ──────────────────────────────────────────
\multicolumn{5}{l}{\textit{NSGA-II + MMDE}} \\[2pt]
$n$
  & Population size
  & \texttt{NSGA\_POP\_SIZE}
  & 100
  & Standard for this scale \\
$G$
  & Number of generations
  & \texttt{NSGA\_GENS}
  & 200
  & Convergence criterion \\
$F$
  & MMDE scaling factor
  & \texttt{MMDE\_F}
  & 0.5
  & Jafari \& Rezvani~\cite{jafari2021} \\
$\mathrm{CR}$
  & MMDE crossover rate
  & \texttt{MMDE\_CR}
  & 0.9
  & Jafari \& Rezvani~\cite{jafari2021} \\
$P_{\text{batch}}$
  & Tasks per offline batch
  & \texttt{NSGA\_BATCH\_SIZE}
  & 100
  & Memory--compute trade-off \\
$N_{\text{batches}}$
  & Offline batches
  & \texttt{N\_OFFLINE\_BATCHES}
  & 1,000
  & Training set size \\

\midrule
%% ── DQN Agent ───────────────────────────────────────────────
\multicolumn{5}{l}{\textit{DQN Agent}} \\[2pt]
$|\mathbf{s}|$
  & State dimension
  & \texttt{AGENT1\_STATE\_DIM}
  & 13
  & Eq.~(21) \\
$|\mathcal{A}|$
  & Action dimension
  & \texttt{AGENT1\_ACTION\_DIM}
  & 5
  & $K+1$ destinations \\
$\bm{h}$
  & Hidden layer widths
  & \texttt{AGENT1\_HIDDEN}
  & [256, 128]
  & Architecture search \\
$\eta$
  & Adam learning rate
  & \texttt{AGENT1\_LR}
  & 0.001
  & Kingma \& Ba (2015) \\
$\gamma$
  & Discount factor
  & \texttt{AGENT1\_GAMMA}
  & 0.95
  & Near-future focused \\
$\varepsilon_0$
  & Initial exploration rate
  & \texttt{AGENT1\_EPSILON\_START}
  & 0.30
  & Standard DQN setting \\
$\varepsilon_{\min}$
  & Minimum exploration rate
  & \texttt{AGENT1\_EPSILON\_END}
  & 0.05
  & Maintain exploration \\
$T_\varepsilon$
  & Epsilon decay horizon (steps)
  & \texttt{AGENT1\_EPSILON\_DECAY}
  & 10,000
  & 10\,min at 10 tasks/s \\
$B$
  & Mini-batch size
  & \texttt{AGENT1\_BATCH\_SIZE}
  & 64
  & GPU memory trade-off \\
$|\mathcal{B}|_{\max}$
  & Replay buffer size
  & \texttt{AGENT1\_BUFFER\_SIZE}
  & 50,000
  & Decorrelation coverage \\
$T_{\text{sync}}$
  & Target network sync (steps)
  & \texttt{AGENT1\_TARGET\_UPDATE}
  & 1,000
  & Stable target period \\
$\delta$
  & Huber loss threshold
  & \texttt{HUBER\_DELTA}
  & 1.0
  & Standard DQN value \\

\midrule
%% ── Reward weights ──────────────────────────────────────────
\multicolumn{5}{l}{\textit{Reward weights}} \\[2pt]
$\omega_L$
  & Latency weight
  & \texttt{REWARD\_WEIGHT\_LATENCY}
  & 0.5
  & Primary objective \\
$\omega_E$
  & Energy weight
  & \texttt{REWARD\_WEIGHT\_ENERGY}
  & 0.3
  & Secondary objective \\
$\omega_V$
  & Violation weight
  & \texttt{REWARD\_WEIGHT\_VIOLATION}
  & 0.2
  & Constraint enforcement \\
$\lambda_{\text{crit}}$
  & Safety multiplier
  & \texttt{DEADLINE\_PENALTY\_CRITICAL}
  & 10.0
  & Safety-critical penalty \\
$E_{\text{ref}}$
  & Energy normalisation ref.\ (J)
  & \texttt{ENERGY\_REF\_JOULES}
  & 0.10
  & Typical pebble energy \\

\midrule
%% ── Mobility ────────────────────────────────────────────────
\multicolumn{5}{l}{\textit{Mobility and handoff}} \\[2pt]
$R_k$
  & Fog coverage radius (m)
  & \texttt{FOG\_COVERAGE\_RADIUS}
  & 250
  & 5G small cell radius \\
$T_{\max}$
  & $T_{\text{exit}}$ normalisation bound (s)
  & \texttt{T\_EXIT\_MAX\_S}
  & 10.0
  & Max useful lookahead \\
$s_{\max}$
  & Speed normalisation bound (m/s)
  & \texttt{SPEED\_MAX\_MS}
  & 33.3
  & 120\,km/h maximum \\

\midrule
%% ── Simulation ──────────────────────────────────────────────
\multicolumn{5}{l}{\textit{Simulation}} \\[2pt]
$N$
  & Number of vehicles
  & \texttt{N\_VEHICLES}
  & 50
  & Urban density scenario \\
$\lambda_{\text{task}}$
  & Task arrival rate per vehicle (Hz)
  & \texttt{TASK\_RATE\_HZ}
  & 10
  & 10\,fps camera \\
$T_{\text{sim}}$
  & Simulation duration (s)
  & \texttt{SIM\_DURATION\_S}
  & 600
  & 10-minute evaluation \\
$T_{\text{warm}}$
  & Warm-up period (s)
  & \texttt{WARMUP\_S}
  & 60
  & Agent stabilisation \\
—
  & Random seed
  & \texttt{RANDOM\_SEED}
  & 42
  & Reproducibility \\
—
  & Independent runs
  & \texttt{N\_RUNS}
  & 5
  & Statistical validity \\

\end{longtable}
\endgroup
% ================================================================
\section{System Model and Notation}
\label{sec:system}
% ================================================================

\subsection{Network Topology}

The system operates over a three-tier hierarchy. Let
$\mathcal{F} = \{f_1, f_2, \ldots, f_K\}$ be the set of $K$ fog nodes and
$\mathcal{C}$ denote the cloud server. The full set of compute destinations is:
\begin{equation}
  \mathcal{X} = \mathcal{F} \cup \{\mathcal{C}\}
\end{equation}

Each fog node $f_k \in \mathcal{F}$ is characterised by a tuple:
\begin{equation}
  f_k = \langle \bm{p}_k,\; R_k,\; \mu_k,\; \rho_k(t) \rangle
\end{equation}
where $\bm{p}_k \in \mathbb{R}^2$ is the fog's two-dimensional position in metres,
$R_k$ is its coverage radius in metres, $\mu_k$ is its compute capacity in
Million Instructions Per Second (MIPS), and $\rho_k(t) \in [0,1]$ is its CPU
utilisation at time $t$, which evolves dynamically as tasks arrive and depart.

The cloud is characterised by:
\begin{equation}
  \mathcal{C} = \langle \mu_c,\; \delta_{\text{WAN}} \rangle
\end{equation}
where $\mu_c > \mu_k$ is the cloud's compute capacity and
$\delta_{\text{WAN}}$ is the fixed wide-area network propagation delay in milliseconds.

\textbf{Static network assumption.}
The network infrastructure — link bandwidths, propagation delays, and fog node
positions — is assumed \emph{static} throughout the system's operational lifetime.
Specifically, the upload bandwidth $B$, the fog-to-cloud backbone bandwidth $B_c$,
and the WAN delay $\delta_{\text{WAN}}$ are fixed constants. This assumption reflects
the real-world constraint that a deployed fog infrastructure does not dynamically
reconfigure its physical links. Only the per-node CPU utilisation $\rho_k(t)$ and
queue depth $q_k(t)$ vary at runtime, driven by the arrival and departure of tasks.
As a consequence, network resource allocation is not an optimisation variable in this
framework; the optimisation concerns exclusively the \emph{compute placement} decision
$\mathbf{x}$ — which node executes each task step.

\subsection{Vehicle Model}

Let $\mathcal{V} = \{v_1, v_2, \ldots, v_N\}$ be a fleet of $N$ mobile IoT
vehicles. At any time $t$, vehicle $v_i$ is described by its \emph{Spatial Tag}:
\begin{equation}
  \sigma_i(t) = \langle \bm{q}_i(t),\; s_i(t),\; \phi_i(t) \rangle
\end{equation}
where $\bm{q}_i(t) \in \mathbb{R}^2$ is the vehicle's GPS-derived position,
$s_i(t)$ is its speed in metres per second, and $\phi_i(t) \in [0, 2\pi)$
is its heading in radians measured from the positive $x$-axis. The Cartesian
velocity vector of vehicle $v_i$ is:
\begin{equation}
  \bm{u}_i(t) = s_i(t) \begin{pmatrix} \cos\phi_i(t) \\ \sin\phi_i(t) \end{pmatrix}
\end{equation}

Vehicle $v_i$ is said to be \emph{in coverage} of fog node $f_k$ at time $t$ if:
\begin{equation}
  \|\bm{q}_i(t) - \bm{p}_k\|_2 \leq R_k
\end{equation}

\subsection{DAG Task Model}

Following the formalization in Zhu et al.~\cite{zhu2025dag}, each vehicle
generates application workloads modeled as Directed Acyclic Graphs. A task
$\tau$ submitted by vehicle $v_i$ at time $t_0$ is defined as:
\begin{equation}
  \tau = \langle G,\; D_\tau,\; v_i,\; t_0 \rangle
\end{equation}
where $G = (\mathcal{V}_G, \mathcal{E}_G)$ is the DAG, $D_\tau$ is the
end-to-end deadline in milliseconds, and $v_i$ is the originating vehicle.

Each node $g_j \in \mathcal{V}_G$ represents one processing step with properties:
\begin{equation}
  g_j = \langle \ell_j,\; d_j^{\text{in}},\; d_j^{\text{out}},\; d_j^{\text{local}} \rangle
\end{equation}
where $\ell_j$ is computational demand in MI, $d_j^{\text{in}}$ and
$d_j^{\text{out}}$ are input and output data sizes in kilobytes, and
$d_j^{\text{local}}$ is the per-step deadline in milliseconds.

A directed edge $(g_j, g_m) \in \mathcal{E}_G$ means step $g_m$ cannot begin
execution until step $g_j$ has completed and transmitted its output. When two
consecutive steps are co-located on the same node, the inter-step transfer time
is zero.

\begin{definition}[Critical Path]
The \emph{critical path} $\Pi \subseteq \mathcal{V}_G$ is the path from the
source step to the sink step that maximises total execution time.
\newline
End-to-end application latency is:
\begin{equation}
  L_\tau(\mathbf{x}) = \sum_{g_j \in \Pi} \Bigl[
    T_{\text{exec}}(g_j, x_j) + T_{\text{tx}}(g_j, g_{j+1}, x_j, x_{j+1})
  \Bigr]
  \label{eq:app_latency}
\end{equation}
where $x_j \in \mathcal{X}$ is the destination assigned to step $g_j$, and
$T_{\text{tx}} = 0$ when $x_j = x_{j+1}$.
\end{definition}

\subsection{Offloading Decision Variable}

The offloading decision for a task $\tau$ with $M$ offloadable steps is a
vector:
\begin{equation}
  \mathbf{x} = (x_1, x_2, \ldots, x_M), \quad x_j \in \mathcal{X}
\end{equation}

Step $g_1$ is always executed on the originating device (it handles local
capture and compression) and is therefore excluded from $\mathbf{x}$.

% ================================================================
\section{Latency and Energy Models}
\label{sec:models}
% ================================================================

\subsection{Transmission Time}

The time to transmit $d$ kilobytes from the IoT device to a fog node over a
channel with nominal bandwidth $B$ Mbps is:
\begin{equation}
  T_{\text{tx}}^{\text{up}}(d, B) = \frac{d \times 8}{B \times 10^3} \times 10^3
  = \frac{8d}{B} \quad \text{[ms]}
  \label{eq:tx_time}
\end{equation}
where the factor of $10^3$ converts seconds to milliseconds. A fixed 5G
access latency $\delta_{5G}$ (one-way, typically 1--3\,ms) is added:
\begin{equation}
  T_{\text{access}}(d) = T_{\text{tx}}^{\text{up}}(d, B) + \delta_{5G}
\end{equation}

For fog-to-cloud transmission with backbone bandwidth $B_c$ and WAN propagation
delay $\delta_{\text{WAN}}$:
\begin{equation}
  T_{\text{tx}}^{\text{cloud}}(d) = \frac{8d}{B_c} + \delta_{\text{WAN}}
  \quad \text{[ms]}
\end{equation}

\subsection{Execution Time on Fog}

The nominal execution time of step $g_j$ on fog node $f_k$ at full capacity is:
\begin{equation}
  T_{\text{exec}}^{0}(g_j, f_k) = \frac{\ell_j}{\mu_k} \times 10^3 \quad \text{[ms]}
\end{equation}

Under non-zero load $\rho_k(t)$, the effective execution time accounting for
CPU time-sharing is:
\begin{equation}
  T_{\text{exec}}(g_j, f_k, t) = \frac{\ell_j}{\mu_k \, (1 - \rho_k(t))}
  \times 10^3 \quad \text{[ms]}
  \label{eq:exec_fog}
\end{equation}

\begin{remark}
Equation~\eqref{eq:exec_fog} assumes a work-conserving processor sharing model.
This is consistent with the model adopted in Jafari and Rezvani~\cite{jafari2021}
and Nouri et al.~\cite{nouri2024tof}. The divisor $(1 - \rho_k(t))$ degrades
toward zero as load approaches unity, capturing queue starvation.
\end{remark}

\subsection{Execution Time on Cloud}

Cloud execution uses dedicated GPU-accelerated servers with fixed capacity $\mu_c$:
\begin{equation}
  T_{\text{exec}}(g_j, \mathcal{C}) = \frac{\ell_j}{\mu_c} \times 10^3
  \quad \text{[ms]}
  \label{eq:exec_cloud}
\end{equation}

Cloud load is assumed manageable due to elastic scaling; no load penalty is applied.

\subsection{Total Step Latency}

The total latency incurred by assigning step $g_j$ to destination $x_j$ is:
\begin{equation}
  L_j(x_j, t) =
  \begin{cases}
    T_{\text{access}}(d_j^{\text{in}}) + T_{\text{exec}}(g_j, x_j, t)
    & \text{if } x_j \in \mathcal{F} \\[6pt]
    T_{\text{tx}}^{\text{cloud}}(d_j^{\text{in}}) + T_{\text{exec}}(g_j, \mathcal{C})
    & \text{if } x_j = \mathcal{C}
  \end{cases}
  \label{eq:step_latency}
\end{equation}

\subsection{Energy Model}

Transmission energy for uploading $d$ kilobytes at transmit power $P_{\text{tx}}$
watts over a link occupying $T_{\text{tx}}$ milliseconds is:
\begin{equation}
  E_{\text{tx}}(d) = P_{\text{tx}} \cdot \frac{8d}{B \times 10^3}
  \quad \text{[Joules]}
  \label{eq:tx_energy}
\end{equation}

Fog computation energy for step $g_j$, assuming a compute energy coefficient
$\kappa$ (Joules per Million Instructions):
\begin{equation}
  E_{\text{comp}}(g_j) = \kappa \cdot \ell_j
  \quad \text{[Joules]}
  \label{eq:comp_energy}
\end{equation}

Total energy cost for assigning step $g_j$ to $x_j$:
\begin{equation}
  E_j(x_j) = E_{\text{tx}}(d_j^{\text{in}})
  + \begin{cases}
    E_{\text{comp}}(g_j) & \text{if } x_j \in \mathcal{F} \\
    \alpha \cdot E_{\text{tx}}(d_j^{\text{in}}) & \text{if } x_j = \mathcal{C}
  \end{cases}
  \label{eq:step_energy}
\end{equation}
where $\alpha > 1$ is a WAN energy penalty factor reflecting the longer
transmission distance to the cloud data centre.

% ================================================================
\section{Multi-Objective Optimisation Problem}
\label{sec:problem}
% ================================================================

\subsection{Problem Formulation}

Let $\mathcal{T} = \{\tau_1, \tau_2, \ldots, \tau_N\}$ be a batch of $N$ tasks,
each with $M$ offloadable steps, giving a joint decision vector
$\mathbf{X} \in \mathcal{X}^{N \times M}$. The offloading problem is formulated as:

\begin{equation}
  \min_{\mathbf{X}} \;\; \bm{f}(\mathbf{X}) =
  \begin{pmatrix} f_1(\mathbf{X}) \\ f_2(\mathbf{X}) \end{pmatrix} =
  \begin{pmatrix}
    \displaystyle\sum_{i=1}^{N} \sum_{j=1}^{M} E_j(x_{ij}) \\[10pt]
    \displaystyle\sum_{i=1}^{N} L_{\tau_i}(\mathbf{x}_i)
  \end{pmatrix}
  \label{eq:biobj}
\end{equation}

subject to:
\begin{align}
  L_{\tau_i}(\mathbf{x}_i) &\leq D_{\tau_i}
    \quad \forall\, i \in \{1,\ldots,N\}
    \label{eq:deadline}\\
  \sum_{i,j:\; x_{ij}=f_k} \ell_{ij} &\leq C_k
    \quad \forall\, f_k \in \mathcal{F}
    \label{eq:capacity}\\
  x_{ij} &\in \mathcal{X}
    \quad \forall\, i, j
    \label{eq:domain}
\end{align}

Constraint~\eqref{eq:deadline} enforces per-task end-to-end deadline satisfaction.
Constraint~\eqref{eq:capacity} bounds aggregate MI load per fog node by its
effective capacity $C_k = \mu_k \cdot (1 - \rho_k^{\min})$, where
$\rho_k^{\min}$ is a minimum reserved utilisation margin.

\begin{proposition}
Problem~\eqref{eq:biobj}--\eqref{eq:domain} is NP-hard. The number of feasible
solutions is $|\mathcal{X}|^{N \times M}$, which grows exponentially with both
the number of tasks and the number of DAG steps. Exact optimisation is
computationally intractable at IoT deployment scale.
\end{proposition}

\begin{definition}[Pareto Dominance]
Solution $\mathbf{X}^{(a)}$ \emph{Pareto-dominates} $\mathbf{X}^{(b)}$,
written $\mathbf{X}^{(a)} \prec \mathbf{X}^{(b)}$, if and only if:
\begin{equation}
  f_k(\mathbf{X}^{(a)}) \leq f_k(\mathbf{X}^{(b)}) \;\; \forall k \in \{1,2\},
  \quad\text{and}\quad
  f_j(\mathbf{X}^{(a)}) < f_j(\mathbf{X}^{(b)}) \;\; \text{for some } j
  \label{eq:dominance}
\end{equation}
\end{definition}

\begin{definition}[Pareto Front]
The \emph{Pareto front} $\mathcal{P}^* \subset \mathcal{X}^{N \times M}$
is the set of all solutions that are not dominated by any other feasible solution:
\begin{equation}
  \mathcal{P}^* = \bigl\{ \mathbf{X} \mid \nexists\, \mathbf{X}' :
  \mathbf{X}' \prec \mathbf{X} \bigr\}
\end{equation}
No element of $\mathcal{P}^*$ can improve on $f_1$ without degrading $f_2$,
and vice versa.
\end{definition}

% ================================================================
\section{Execution Cost Classification}
\label{sec:tof}
% ================================================================

\subsection{Execution Cost Metric}

The Execution Cost (EC) of a step $g_j$, introduced by Nouri et al.~\cite{nouri2024tof},
is the time the step would occupy a fog node running at full capacity:
\begin{equation}
  \mathrm{EC}(g_j) = \frac{\ell_j}{\mu_k} \quad \text{[seconds]}
  \label{eq:ec}
\end{equation}

Since $\mu_k$ is uniform across fog nodes in this deployment, EC is
task-dependent only and can be computed without knowledge of which fog node
will be assigned.

\subsection{Boulder and Pebble Classification}

Let $\theta > 0$ be the EC threshold parameter. Step $g_j$ is classified as:
\begin{equation}
  \mathrm{class}(g_j) =
  \begin{cases}
    \textsc{Boulder} & \text{if } \mathrm{EC}(g_j) \geq \theta \\
    \textsc{Pebble}  & \text{if } \mathrm{EC}(g_j) < \theta
  \end{cases}
  \label{eq:classification}
\end{equation}

\textsc{Boulders} are immediately routed to cloud without entering the
optimisation pipeline. \textsc{Pebbles} are queued for multi-objective
optimisation. This partitioning reduces the search space of
problem~\eqref{eq:biobj} by excluding steps that cannot feasibly complete on
any fog node within the deadline.

\subsection{Aggregator and Avalanche Prevention}

Let $Q(t)$ denote the pebble queue depth at time $t$ and $Q_{\max}$ a
configurable overflow threshold. The \emph{Aggregator} activates when:
\begin{equation}
  Q(t) > Q_{\max}
  \label{eq:aggregator_trigger}
\end{equation}

Upon activation, all queued pebbles are bundled into a \emph{super-task}
$\mathcal{S}$ and routed to cloud:
\begin{equation}
  \mathcal{S} = \bigl\{ g_j \in Q(t) \;\big|\; Q(t) > Q_{\max} \bigr\}
\end{equation}

This prevents queue cascade failure (the Pebble Avalanche) identified as an
open gap in the literature~\cite{nouri2024tof}. The total latency of a
super-task is bounded by:
\begin{equation}
  L_{\mathcal{S}} \leq \frac{8 \cdot \sum_{g_j \in \mathcal{S}} d_j^{\text{in}}}{B_c}
  + \delta_{\text{WAN}}
  + \max_{g_j \in \mathcal{S}} T_{\text{exec}}(g_j, \mathcal{C})
\end{equation}
assuming parallel cloud execution of all bundled steps.

% ================================================================
\section{MMDE-Enhanced NSGA-II}
\label{sec:nsga2}
% ================================================================

\subsection{Chromosome Representation}

A candidate solution (chromosome) for a batch of $P$ pebble steps is a
vector of integer genes:
\begin{equation}
  \mathbf{c} = (c_1, c_2, \ldots, c_P), \quad
  c_j \in \{0, 1, \ldots, |\mathcal{F}|\}
  \label{eq:chromosome}
\end{equation}
where gene value $c_j = k$ assigns step $g_j$ to fog node $f_k$, and
$c_j = |\mathcal{F}|$ assigns it to the cloud. The population at generation
$r$ is $\mathbf{C}^{(r)} = \{\mathbf{c}^{(r)}_1, \ldots, \mathbf{c}^{(r)}_n\}$
where $n = |\mathbf{C}|$ is the population size.

\subsection{Objective Evaluation}

Each chromosome $\mathbf{c}$ is evaluated on both objectives:
\begin{align}
  F_1(\mathbf{c}) &= \sum_{j=1}^{P} E_j(x_j(\mathbf{c}))
  \label{eq:obj1}\\
  F_2(\mathbf{c}) &= \sum_{j=1}^{P} L_j(x_j(\mathbf{c}), t)
  \label{eq:obj2}
\end{align}
where $x_j(\mathbf{c})$ maps gene $c_j$ to its corresponding destination in
$\mathcal{X}$.

\subsection{Non-Dominated Sorting}

NSGA-II partitions the population into Pareto fronts $\mathcal{F}_1, \mathcal{F}_2,
\ldots$ via non-dominated sorting~\cite{jafari2021}. For each chromosome
$\mathbf{c}^{(r)}_i$, its domination count $n_i$ (number of solutions that
dominate it) and its dominated set $S_i$ (solutions it dominates) are computed.
Front $\mathcal{F}_1$ contains all solutions with $n_i = 0$. Subsequent fronts
are obtained by successively removing $\mathcal{F}_1$ and recomputing.

\subsection{Crowding Distance}

Within each front, crowding distance $d_i^{\text{crowd}}$ measures the
perimeter of the cuboid formed by the nearest neighbours of solution $i$
in the objective space. For two objectives:
\begin{equation}
  d_i^{\text{crowd}} = \sum_{k=1}^{2}
  \frac{F_k(\mathbf{c}_{i+1}) - F_k(\mathbf{c}_{i-1})}
  {F_k^{\max} - F_k^{\min}}
  \label{eq:crowding}
\end{equation}
where solutions are sorted by objective $k$ and boundary solutions receive
$d_i^{\text{crowd}} = \infty$. Larger crowding distance is preferred when
two solutions share the same front rank, ensuring diversity preservation.

\subsection{Minimax Differential Evolution Mutation}

Standard NSGA-II uses polynomial mutation, which samples perturbations
uniformly and wastes evaluations on unproductive regions. The Minimax
Differential Evolution (MMDE) operator, proposed in~\cite{jafari2021},
replaces this with a directional mutation derived from population structure.

Given three distinct solutions $\mathbf{r}_1, \mathbf{r}_2, \mathbf{r}_3
\in \mathbf{C}^{(r)}$, the mutant vector is:
\begin{equation}
  \mathbf{V} = \mathbf{r}_1 + F \cdot (\mathbf{r}_2 - \mathbf{r}_3)
  \label{eq:mmde}
\end{equation}
where $F \in (0, 2)$ is the differential scaling factor. For integer gene
representations, the mutation of gene $j$ is:
\begin{equation}
  V_j = \mathrm{clip}\bigl(
    r_{1,j} + \lfloor F \cdot (r_{2,j} - r_{3,j}) + 0.5 \rfloor,
    \; 0,\; |\mathcal{X}| - 1
  \bigr)
  \label{eq:mmde_int}
\end{equation}

Binomial crossover then combines the mutant with the target chromosome
$\mathbf{c}_i^{(r)}$ to produce a trial chromosome $\mathbf{u}_i$:
\begin{equation}
  u_{i,j} =
  \begin{cases}
    V_j              & \text{if } \xi_j \leq \mathrm{CR} \text{ or } j = j_{\text{rand}} \\
    c_{i,j}^{(r)}   & \text{otherwise}
  \end{cases}
  \label{eq:crossover}
\end{equation}
where $\xi_j \sim \mathcal{U}(0,1)$ is drawn independently per gene,
$\mathrm{CR} \in (0,1)$ is the crossover rate, and $j_{\text{rand}}$ is a
randomly chosen index ensuring at least one gene is taken from the mutant.

The trial chromosome $\mathbf{u}_i$ replaces $\mathbf{c}_i^{(r)}$ in the
next generation if it is not dominated:
\begin{equation}
  \mathbf{c}_i^{(r+1)} =
  \begin{cases}
    \mathbf{u}_i       & \text{if } \mathbf{u}_i \prec \mathbf{c}_i^{(r)}
                         \text{ or } \mathbf{u}_i = \mathbf{c}_i^{(r)} \\
    \mathbf{c}_i^{(r)} & \text{otherwise}
  \end{cases}
  \label{eq:selection}
\end{equation}

The full MMDE-NSGA-II algorithm is summarised in Algorithm~\ref{alg:nsga2}.

\begin{algorithm}[H]
\caption{TOF-MMDE-NSGA-II Offline Optimiser}
\label{alg:nsga2}
\begin{algorithmic}[1]
\State \textbf{Input:} pebble step set $\mathcal{G}_P$, fog state
  $\{\rho_k\}$, population size $n$, generations $G$,
  scaling factor $F$, crossover rate $\mathrm{CR}$
\State \textbf{Output:} Pareto front $\mathcal{P}^*$, knee point $\mathbf{c}^*$
\State Initialise population $\mathbf{C}^{(0)}$ with $n$ random chromosomes
\State Evaluate $\bm{F}(\mathbf{c})$ for all $\mathbf{c} \in \mathbf{C}^{(0)}$
\State Compute non-dominated fronts and crowding distances
\For{$r = 1$ \textbf{to} $G$}
  \For{each $\mathbf{c}_i^{(r-1)} \in \mathbf{C}^{(r-1)}$}
    \State Sample $\mathbf{r}_1, \mathbf{r}_2, \mathbf{r}_3$ from $\mathbf{C}^{(r-1)}$, all distinct
    \State Compute mutant $\mathbf{V}$ via Eq.~\eqref{eq:mmde_int}
    \State Compute trial $\mathbf{u}_i$ via Eq.~\eqref{eq:crossover}
    \State Evaluate $\bm{F}(\mathbf{u}_i)$
    \State Select $\mathbf{c}_i^{(r)}$ via Eq.~\eqref{eq:selection}
  \EndFor
  \State Combine $\mathbf{C}^{(r)} \cup \{$offspring$\}$, re-sort, select best $n$
\EndFor
\State Extract Pareto front $\mathcal{P}^* = \mathcal{F}_1$ from final population
\State Compute knee point $\mathbf{c}^*$ via Eq.~\eqref{eq:knee}
\State \Return $\mathcal{P}^*, \mathbf{c}^*$
\end{algorithmic}
\end{algorithm}

\subsection{Knee Point Selection}

Let $\bm{F}^{\text{utopia}} = (F_1^{\min}, F_2^{\min})$ be the utopia point,
where each component is the minimum of that objective over the Pareto front
(achievable only if the objectives were independent). The knee point
$\mathbf{c}^*$ is the Pareto-front solution with minimum normalised Euclidean
distance to the utopia point:
\begin{equation}
  \mathbf{c}^* = \arg\min_{\mathbf{c} \in \mathcal{P}^*}
  \left\|
  \begin{pmatrix}
    \dfrac{F_1(\mathbf{c}) - F_1^{\min}}{F_1^{\max} - F_1^{\min}} \\[10pt]
    \dfrac{F_2(\mathbf{c}) - F_2^{\min}}{F_2^{\max} - F_2^{\min}}
  \end{pmatrix}
  \right\|_2
  \label{eq:knee}
\end{equation}

The knee point provides the best balance between energy minimisation and
latency minimisation without requiring any preference weights to be specified
a priori.

% ================================================================
\section{Deep Q-Network Task Placement Agent}
\label{sec:agent}
% ================================================================

\subsection{Markov Decision Process Formulation}

The real-time task placement problem is modeled as a Markov Decision Process
(MDP) $\mathcal{M} = (\mathcal{S}, \mathcal{A}, \mathcal{R}, \mathcal{P},
\gamma)$, where:

\begin{itemize}[leftmargin=2em]
  \item $\mathcal{S}$ is the state space — the agent's observable view of the
    system at the moment a pebble step arrives for placement.
  \item $\mathcal{A} = \mathcal{X} = \{f_1, f_2, \ldots, f_K, \mathcal{C}\}$
    is the action space — the destination to which the step is assigned.
  \item $\mathcal{R}: \mathcal{S} \times \mathcal{A} \to \mathbb{R}$ is the
    reward function.
  \item $\mathcal{P}(s' \mid s, a)$ is the transition probability — here
    implicit in the simulation dynamics.
  \item $\gamma \in (0,1)$ is the discount factor.
\end{itemize}

\subsection{State Representation}

Under the static network assumption, the upload bandwidth $B$ and
all link capacities are fixed constants known at design time. The agent therefore
does not need to observe network utilisation as a runtime state variable. The state
vector $\mathbf{s} \in \mathbb{R}^{11}$ observed by the agent at decision time $t$
for step $g_j$ submitted by vehicle $v_i$ is:
\begin{equation}
  \mathbf{s} = \bigl(
    \rho_1(t),\; \rho_2(t),\; \rho_3(t),\; \rho_4(t),\;
    q_1(t),\; q_2(t),\; q_3(t),\; q_4(t),\;
    \widehat{\mathrm{EC}}_j,\;
    \widehat{s}_i(t),\; \widehat{T}_{\text{exit}}
  \bigr)
  \label{eq:state}
\end{equation}

where hats denote normalisation to $[0,1]$. The components are:
\begin{itemize}[leftmargin=2em]
  \item $\rho_k(t) \in [0,1]$: CPU utilisation of fog node $f_k$ ($k = 1,\ldots,4$).
    This is the primary signal for load-aware placement.
  \item $q_k(t) / q_{\max}$: queue depth of fog node $f_k$ normalised by the maximum
    queue depth $q_{\max}$. Captures queuing latency not reflected in CPU utilisation alone.
  \item $\widehat{\mathrm{EC}}_j = \min(\mathrm{EC}(g_j) / \theta, 1)$: normalised
    Execution Cost of the current step. Heavier steps require lighter-loaded nodes.
  \item $\widehat{s}_i(t) = s_i(t) / s_{\max}$: vehicle speed normalised by $s_{\max}$.
    Faster vehicles exit fog zones sooner, affecting placement urgency.
  \item $\widehat{T}_{\text{exit}} = \min(T_{\text{exit}} / T_{\max}, 1)$: normalised
    time until the vehicle exits the current fog zone. The core proactive handoff signal.
\end{itemize}

\begin{remark}
Two dimensions present in earlier formulations --- upload link utilisation
$\widehat{B}(t)$ and remaining deadline fraction $\widehat{d}_j^{\text{rem}}$ ---
are removed. Under the static network assumption, $B$ is a fixed constant and
contributes no information for placement decisions. The deadline fraction is
implicitly captured through the $T_{\text{exit}}$ dimension: if a vehicle is about
to leave the zone, the effective deadline for fog completion is already constrained.
Reducing the state dimension from 13 to 11 shrinks the input weight matrix
$\bm{W}^{(1)}$ from $\mathbb{R}^{256 \times 13}$ to $\mathbb{R}^{256 \times 11}$,
marginally improving training efficiency without sacrificing decision quality.
\end{remark}

\subsection{Q-Network Architecture}

The action-value function $Q(s, a;\, \bm{\theta})$ is approximated by a
fully-connected neural network with parameters $\bm{\theta}$:
\begin{align}
  \bm{h}^{(1)} &= \mathrm{ReLU}\!\left(\bm{W}^{(1)} \mathbf{s} + \bm{b}^{(1)}\right),
    \quad \bm{W}^{(1)} \in \mathbb{R}^{256 \times 13}
    \label{eq:layer1}\\
  \bm{h}^{(2)} &= \mathrm{ReLU}\!\left(\bm{W}^{(2)} \bm{h}^{(1)} + \bm{b}^{(2)}\right),
    \quad \bm{W}^{(2)} \in \mathbb{R}^{128 \times 256}
    \label{eq:layer2}\\
  \bm{Q}(s;\,\bm{\theta}) &= \bm{W}^{(3)} \bm{h}^{(2)} + \bm{b}^{(3)},
    \quad \bm{W}^{(3)} \in \mathbb{R}^{|\mathcal{A}| \times 128}
    \label{eq:output}
\end{align}

The agent selects actions via $\varepsilon$-greedy policy:
\begin{equation}
  a_t =
  \begin{cases}
    \text{random } a \in \mathcal{A} & \text{with probability } \varepsilon_t \\
    \arg\max_{a \in \mathcal{A}} Q(s_t, a;\, \bm{\theta}) & \text{otherwise}
  \end{cases}
  \label{eq:epsilon_greedy}
\end{equation}

Exploration rate $\varepsilon_t$ decays linearly:
\begin{equation}
  \varepsilon_t = \max\!\left(
    \varepsilon_{\min},\;
    \varepsilon_0 - \frac{(\varepsilon_0 - \varepsilon_{\min})}{T_{\varepsilon}} \cdot t
  \right)
  \label{eq:epsilon_decay}
\end{equation}
where $\varepsilon_0$ is the initial exploration rate, $\varepsilon_{\min}$
is the final minimum, and $T_\varepsilon$ is the decay horizon in steps.

\subsection{Reward Function}

The reward received after step $g_j$ completes with latency $L_j$,
energy $E_j$, and deadline $d_j^{\text{local}}$ is:
\begin{equation}
  R(g_j, x_j) = -\omega_L \cdot \widetilde{L}_j
    - \omega_E \cdot \widetilde{E}_j
    - \omega_V \cdot \mathbf{1}[L_j > d_j^{\text{local}}] \cdot \lambda
  \label{eq:reward}
\end{equation}

where:
\begin{align}
  \widetilde{L}_j &= \min\!\left(\frac{L_j}{d_j^{\text{local}}},\; 3\right)
    \label{eq:norm_lat}\\
  \widetilde{E}_j &= \min\!\left(\frac{E_j}{E_{\text{ref}}},\; 3\right)
    \label{eq:norm_energy}\\
  \mathbf{1}[\cdot] &\in \{0, 1\} \text{ is the deadline violation indicator}
    \label{eq:indicator}\\
  \lambda &\in \{1, \lambda_{\text{crit}}\}
    \text{ is the safety multiplier}
    \label{eq:lambda}
\end{align}

$\omega_L, \omega_E, \omega_V > 0$ are reward weights satisfying
$\omega_L + \omega_E + \omega_V = 1$.
$\lambda_{\text{crit}} \gg 1$ applies to safety-critical steps such as
collision avoidance alerts, ensuring the agent prioritises their deadlines.

\subsection{Experience Replay and Temporal Difference Learning}

Transitions $(s_t, a_t, r_t, s_{t+1}, \text{done}_t)$ are stored in a replay
buffer $\mathcal{B}$ of capacity $|\mathcal{B}|_{\max}$. At each update step,
a mini-batch $\mathcal{B}' \subset \mathcal{B}$ of size $B$ is sampled uniformly.
The TD target for each transition in $\mathcal{B}'$ is computed using a
target network $\bm{\theta}^{-}$, which is a periodically-frozen copy of
the online network:
\begin{equation}
  y_t = r_t + \gamma \cdot
  \max_{a' \in \mathcal{A}} Q(s_{t+1}, a';\, \bm{\theta}^{-})
  \cdot (1 - \text{done}_t)
  \label{eq:td_target}
\end{equation}

The online network is updated by minimising the Huber loss over the mini-batch:
\begin{equation}
  \mathcal{L}(\bm{\theta}) = \frac{1}{B} \sum_{(s,a,r,s') \in \mathcal{B}'}
  \mathcal{H}_\delta\!\left(y - Q(s, a;\, \bm{\theta})\right)
  \label{eq:huber}
\end{equation}
\begin{equation}
  \mathcal{H}_\delta(u) =
  \begin{cases}
    \tfrac{1}{2} u^2              & \text{if } |u| \leq \delta \\
    \delta (|u| - \tfrac{1}{2}\delta) & \text{otherwise}
  \end{cases}
  \label{eq:huber_def}
\end{equation}

Gradients are computed by back-propagation and applied via Adam optimiser
with learning rate $\eta$. The target network is synchronised every
$T_{\text{sync}}$ update steps:
\begin{equation}
  \bm{\theta}^{-} \leftarrow \bm{\theta}
  \quad \text{every } T_{\text{sync}} \text{ steps}
  \label{eq:target_sync}
\end{equation}

\subsection{Behavioral Cloning Pre-training}

To eliminate cold-start degradation, the agent is pre-trained from NSGA-II
Pareto solutions before live deployment. From each offline optimisation run,
the knee-point chromosome $\mathbf{c}^*$ is decomposed into supervised
training pairs:
\begin{equation}
  \mathcal{D}_{\text{BC}} = \bigl\{
    (\mathbf{s}_j,\; a_j^*) \;\big|\;
    a_j^* = x_j(\mathbf{c}^*),\; j = 1,\ldots,P
  \bigr\}
  \label{eq:bc_dataset}
\end{equation}

where $\mathbf{s}_j$ is the state vector at the moment step $g_j$ would be
placed, and $a_j^*$ is the NSGA-II optimal action for that step. The network
is trained on $\mathcal{D}_{\text{BC}}$ by minimising cross-entropy loss:
\begin{equation}
  \mathcal{L}_{\text{BC}}(\bm{\theta}) =
  -\frac{1}{|\mathcal{D}_{\text{BC}}|}
  \sum_{(\mathbf{s}, a^*) \in \mathcal{D}_{\text{BC}}}
  \log \frac{\exp Q(\mathbf{s}, a^*;\, \bm{\theta})}
  {\sum_{a \in \mathcal{A}} \exp Q(\mathbf{s}, a;\, \bm{\theta})}
  \label{eq:bc_loss}
\end{equation}

This initialises the agent's policy close to NSGA-II's optimal policy,
eliminating the initial period of near-random decisions.

\subsection{Offline Training Data Generation}

The offline training dataset $\mathcal{D}_{\text{BC}}$ is constructed by
running TOF-MMDE-NSGA-II on $N_{\text{batches}}$ synthetic batches, each
representing a plausible historical system state. For each batch $b \in
\{1,\ldots,N_{\text{batches}}\}$, the fog state is sampled as follows:

\begin{equation}
  \rho_k^{(b)} \sim \mathcal{U}(0.20,\; 0.75),
  \quad k = 1,\ldots,K
  \label{eq:load_sample}
\end{equation}
\begin{equation}
  q_k^{(b)} \sim \mathcal{U}(0,\; Q_{\max}),
  \quad k = 1,\ldots,K
  \label{eq:queue_sample}
\end{equation}
\begin{equation}
  s_i^{(b)} \sim \mathcal{N}(\mu_s,\; \sigma_s^2),
  \quad \phi_i^{(b)} \sim \mathcal{U}(0,\; 2\pi)
  \label{eq:vehicle_sample}
\end{equation}

where $\mu_s = 60/3.6$\,m/s and $\sigma_s = 15/3.6$\,m/s are the mean and
standard deviation of vehicle speed. Task MI values for pebble steps are drawn
from the empirical distribution of the application DAG:
\begin{equation}
  \ell_j^{(b)} \sim p_{\text{MI}},
  \quad p_{\text{MI}} = \{200, 500, 800, 1000, 1200, 1500, 1800\}
  \text{ with equal probability}
  \label{eq:mi_sample}
\end{equation}

Each batch yields $P_{\text{batch}}$ pebble steps. After running NSGA-II and
selecting the knee point $\mathbf{c}^{*(b)}$, the batch contributes
$P_{\text{batch}}$ labelled training pairs. The full dataset is:
\begin{equation}
  \mathcal{D}_{\text{BC}} = \bigcup_{b=1}^{N_{\text{batches}}}
  \bigl\{(\mathbf{s}_j^{(b)},\; a_j^{*(b)}) \;\big|\; j=1,\ldots,P_{\text{batch}}\bigr\}
  \label{eq:full_dataset}
\end{equation}
with total size $|\mathcal{D}_{\text{BC}}| = N_{\text{batches}} \cdot P_{\text{batch}}
= 10^5$ labelled examples at default settings.

\subsection{Offline-to-Online Transition}

Pre-training terminates when either the behavioral cloning loss falls below
a convergence threshold $\epsilon_{\text{BC}}$ or the maximum number of
epochs $E_{\text{BC}}$ is reached:
\begin{equation}
  \text{stop pre-training when: }
  \mathcal{L}_{\text{BC}}(\bm{\theta}) < \epsilon_{\text{BC}}
  \quad \text{or} \quad
  e = E_{\text{BC}}
  \label{eq:bc_stop}
\end{equation}

The weights at termination define the initial online policy:
\begin{equation}
  \bm{\theta}_0 = \arg\min_{e \leq E_{\text{BC}}}
  \mathcal{L}_{\text{BC}}(\bm{\theta}^{(e)})
  \label{eq:theta_init}
\end{equation}

Both the online network and the target network are initialised to
$\bm{\theta}_0$ before live simulation begins:
\begin{equation}
  \bm{\theta} \leftarrow \bm{\theta}_0,
  \quad \bm{\theta}^{-} \leftarrow \bm{\theta}_0
  \label{eq:init_both}
\end{equation}

From this point, the agent operates under the $\varepsilon$-greedy policy
of Eq.~\eqref{eq:epsilon_greedy} and updates its weights via TD learning
as each task step completes. The exploration rate $\varepsilon$ is reset to
$\varepsilon_0$ at the start of online deployment, ensuring continued
exploration despite the informed initialisation.


% ================================================================
\section{Proactive Mobility Management}
\label{sec:mobility}
% ================================================================

\subsection{Zone Exit Time Estimation}

The proactive handoff mechanism operates by comparing the estimated
execution time of a step against the time available before the vehicle
exits the serving fog node's coverage zone.

\begin{definition}[Closing Speed]
The closing speed $v_{\text{close}}(v_i, f_k, t)$ of vehicle $v_i$ with
respect to fog node $f_k$ at time $t$ is the radial component of the
vehicle's velocity directed outward from the fog centre:
\begin{equation}
  v_{\text{close}} =
  \bm{u}_i(t) \cdot \hat{\bm{n}}_{ik}(t)
  \label{eq:v_close}
\end{equation}
where $\hat{\bm{n}}_{ik}(t) = -(\bm{p}_k - \bm{q}_i(t)) /
\|\bm{p}_k - \bm{q}_i(t)\|_2$ is the unit outward radial vector
from the fog centre toward the vehicle.
\end{definition}

\begin{definition}[Zone Exit Time]
The \emph{zone exit time} $T_{\text{exit}}(v_i, f_k, t)$ is the time in
seconds until vehicle $v_i$ crosses the boundary of fog node $f_k$'s
coverage zone, under linear extrapolation of current velocity:
\begin{equation}
  T_{\text{exit}}(v_i, f_k, t) =
  \begin{cases}
    \dfrac{R_k - \|\bm{q}_i(t) - \bm{p}_k\|_2}{v_{\text{close}}}
    & \text{if } v_{\text{close}} > 0 \\[8pt]
    +\infty & \text{if } v_{\text{close}} \leq 0
  \end{cases}
  \label{eq:texit}
\end{equation}
$v_{\text{close}} \leq 0$ means the vehicle is moving toward or parallel
to the fog centre, so no imminent exit is predicted.
\end{definition}

\subsection{Handoff Mode Selection}

For each pebble step $g_j$ assigned by the DQN agent to fog node $f_k$,
the mobility management layer computes both $T_{\text{exit}}$ and the
effective execution time $T_{\text{exec}}(g_j, f_k, t)$ via
Eq.~\eqref{eq:exec_fog}. The handoff mode is then selected as:
\begin{equation}
  \mathrm{mode}(g_j, f_k, t) =
  \begin{cases}
    \textsc{Direct}     & \text{if } T_{\text{exec}} < T_{\text{exit}} \\
    \textsc{Proactive}  & \text{if } T_{\text{exec}} \geq T_{\text{exit}}
  \end{cases}
  \label{eq:mode}
\end{equation}

\textbf{\textsc{Direct} mode:} The step will complete before the vehicle exits.
No migration action is taken.

\textbf{\textsc{Proactive} mode:} The step will not complete before the vehicle
exits. The predicted destination fog node $f_{k'}$ is identified by projecting
the vehicle's position $T_{\text{exit}}$ seconds forward:
\begin{equation}
  \bm{q}_i^{\text{pred}} = \bm{q}_i(t) + T_{\text{exit}} \cdot \bm{u}_i(t)
  \label{eq:position_pred}
\end{equation}
The destination node is:
\begin{equation}
  f_{k'} = \arg\min_{f_k \in \mathcal{F} \setminus \{f_k\}}
  \bigl\|\bm{q}_i^{\text{pred}} - \bm{p}_k\bigr\|_2
  \quad \text{s.t.} \quad
  \bigl\|\bm{q}_i^{\text{pred}} - \bm{p}_{k'}\bigr\|_2 \leq R_{k'}
  \label{eq:dest_node}
\end{equation}
A container pre-spin request is immediately issued to $f_{k'}$, reserving
$\mu_{k'} \cdot T_{\text{exec}}(g_j, f_{k'}, t)$ compute time and
$d_j^{\text{in}}$ kilobytes of memory for the arriving step. The vehicle
connects to $f_{k'}$ at exit time and the step continues without interruption.

\subsection{Reactive Fallback: NTB/HTB System}

When trajectory prediction fails (unexpected turn, GPS signal loss), the
reactive fallback activates. Tasks in the Normal Task Buffer (NTB) are
processed in arrival order. Upon detecting vehicle disconnection, in-flight
steps are moved to the Handoff Task Buffer (HTB):
\begin{equation}
  \mathrm{HTB} \leftarrow \mathrm{HTB} \cup \{g_j \mid g_j \in \mathrm{NTB},\;
  v_i(g_j) \text{ disconnected}\}
  \label{eq:htb}
\end{equation}

The HTB is assigned strict priority:
\begin{equation}
  \text{scheduled CPU fraction for HTB} = 1.0
  \quad \text{while } |\mathrm{HTB}| > 0
  \label{eq:htb_priority}
\end{equation}

Completed results are held in a Central Tracker indexed by vehicle ID.
Upon reconnection of vehicle $v_i$ to any fog node $f_k$:
\begin{equation}
  \text{deliver all results } \in \mathrm{Tracker}[v_i]
  \text{ to } v_i \text{ via } f_k
  \label{eq:delivery}
\end{equation}
This eliminates task re-submission and the associated retransmission cost.

\subsection{Step-Level Assessment for DAG Tasks}

Prior work assesses $T_{\text{exit}}$ once per task~\cite{mavto2024}.
This framework evaluates Eq.~\eqref{eq:texit} independently for every
DAG step $g_j$, against the specific fog node $x_j$ assigned by the DQN agent:
\begin{equation}
  \forall g_j \in \Pi: \quad
  \mathrm{mode}(g_j) = \mathrm{mode}(g_j, x_j(\mathbf{c}^*), t_j)
  \label{eq:per_step}
\end{equation}
where $t_j$ is the time at which step $g_j$ begins transmission. Since
different steps are assigned to different fog nodes, and the vehicle position
evolves as each step executes, this per-step assessment produces more
accurate migration decisions than task-level assessment.

% ================================================================
% ================================================================
\section{Computational Complexity Analysis}
\label{sec:complexity}
% ================================================================

\subsection{Offline Phase: TOF-MMDE-NSGA-II}

The dominant cost in the offline phase is NSGA-II. Let $n$ be the population
size, $G$ the number of generations, $P$ the number of pebble steps per batch,
and $K$ the number of fog nodes. The cost per generation is:

\begin{equation}
  \mathcal{O}_{\text{gen}} = O(n^2) + O(n \cdot P \cdot K)
  \label{eq:nsga_gen_cost}
\end{equation}

The $O(n^2)$ term accounts for non-dominated sorting, which compares every
pair of solutions~\cite{deb2002nsga2}. The $O(n \cdot P \cdot K)$ term
accounts for objective evaluation: for each of the $n$ chromosomes, all
$P$ genes are evaluated against all $K$ destinations. Total offline cost
across $G$ generations and $N_{\text{batches}}$ batches is:
\begin{equation}
  \mathcal{O}_{\text{offline}} =
  N_{\text{batches}} \cdot G \cdot \bigl(O(n^2) + O(n \cdot P \cdot K)\bigr)
  \label{eq:offline_total}
\end{equation}

At default settings ($N_{\text{batches}} = 10^3$, $G = 200$, $n = 100$,
$P = 100$, $K = 4$), the offline phase runs once before deployment and its
cost is amortised over the full simulation lifetime.

\subsection{Online Phase: DQN Inference}

At runtime, the agent makes one placement decision per pebble step via a
single forward pass through the Q-network. The cost of one inference is:
\begin{equation}
  \mathcal{O}_{\text{infer}} =
  O(|\mathbf{s}| \cdot h_1 + h_1 \cdot h_2 + h_2 \cdot |\mathcal{A}|)
  \label{eq:dqn_infer}
\end{equation}

With $|\mathbf{s}| = 11$, $h_1 = 256$, $h_2 = 128$, $|\mathcal{A}| = 5$:
\begin{equation}
  \mathcal{O}_{\text{infer}} = O(11 \cdot 256 + 256 \cdot 128 + 128 \cdot 5)
  = O(35,\!776)
  \label{eq:dqn_infer_numeric}
\end{equation}

This is a fixed constant independent of the number of tasks, fog nodes, or
vehicles. On modern hardware a single inference executes in under 1\,ms,
satisfying the real-time constraint imposed by the 200\,ms end-to-end deadline.
This is the key computational advantage over running NSGA-II at inference time,
which would require $O(n \cdot P \cdot K \cdot G)$ operations per decision.

\subsection{Mobility Prediction}

Computing $T_{\text{exit}}$ via Eq.~\eqref{eq:texit} requires one dot product
and one division: $O(2)$ per step. The proactive destination search
Eq.~\eqref{eq:dest_node} evaluates $K - 1$ distance comparisons: $O(K)$.
The total per-step mobility cost is $O(K)$, which at $K = 4$ is negligible.

\subsection{Summary}

\begin{center}
\renewcommand{\arraystretch}{1.3}
\begin{tabular}{@{}lll@{}}
\toprule
\textbf{Component} & \textbf{When} & \textbf{Complexity per call} \\
\midrule
TOF-Broker (EC classification) & per step & $O(1)$ \\
NSGA-II per generation         & offline  & $O(n^2 + n \cdot P \cdot K)$ \\
Knee-point selection           & offline  & $O(n)$ \\
DQN inference (Agent 1)        & per step & $O(36{,}288)$ constant \\
DQN weight update              & per step & $O(B \cdot h_1 \cdot h_2)$ \\
$T_{\text{exit}}$ computation  & per step & $O(K)$ \\
HTB queue operation            & on event & $O(1)$ \\
\bottomrule
\end{tabular}
\end{center}

% \begin{table}[H]
% \centering
% \caption{Complete parameter set of the PCNME framework.}
% \label{tab:params}
% \renewcommand{\arraystretch}{1.25}
% \begin{tabularx}{\textwidth}{@{}L{1.8cm}L{5.2cm}C{1.5cm}L{4.5cm}@{}}
% \toprule
% \textbf{Symbol} & \textbf{Description} & \textbf{Default} & \textbf{Source / Justification} \\
% \midrule
% \multicolumn{4}{l}{\textit{Compute and network}} \\
% $\mu_k$     & Fog node MIPS capacity          & 2,000        & Edge server benchmark \\
% $\mu_c$     & Cloud MIPS capacity             & 8,000        & 4$\times$ fog (GPU server) \\
% $B$         & IoT upload bandwidth (Mbps)     & 100          & 5G NR sub-6 GHz typical \\
% $B_c$       & Fog-cloud backbone (Mbps)       & 1,000        & Urban fibre link \\
% $\delta_{5G}$   & 5G one-way latency (ms)     & 2            & 3GPP Release 16 \\
% $\delta_{\text{WAN}}$ & Fog-cloud WAN delay (ms) & 30        & Istanbul--cloud typical \\
% $\alpha$    & WAN energy penalty factor       & 1.8          & Longer path coefficient \\
% $\kappa$    & Fog compute energy (J/MI)       & 0.001        & Calibrated to ARM server \\
% $P_{\text{tx}}$ & Device transmit power (W)   & 0.5          & 5G UE typical \\
% \midrule
% \multicolumn{4}{l}{\textit{TOF-Broker}} \\
% $\theta$    & EC classification threshold (s) & 1.0          & \cite{nouri2024tof} \\
% $Q_{\max}$ & Aggregator trigger depth        & 50           & Empirically tuned \\
% \midrule
% \multicolumn{4}{l}{\textit{NSGA-II + MMDE}} \\
% $n$         & Population size                 & 100          & Standard for this scale \\
% $G$         & Number of generations           & 200          & Convergence criterion \\
% $F$         & MMDE scaling factor             & 0.5          & \cite{jafari2021} \\
% $\mathrm{CR}$ & MMDE crossover rate          & 0.9          & \cite{jafari2021} \\
% $P_{\text{batch}}$ & Tasks per offline batch  & 100          & Memory-compute trade-off \\
% $N_{\text{batches}}$ & Offline batches        & 1,000        & Training set size \\
% \midrule
% \multicolumn{4}{l}{\textit{DQN Agent}} \\
% $|\mathbf{s}|$  & State dimension             & 13           & Eq.~\eqref{eq:state} \\
% $|\mathcal{A}|$ & Action dimension            & 5            & $K + 1$ destinations \\
% $\bm{h}$    & Hidden layer widths             & [256, 128]   & Architecture search \\
% $\eta$      & Adam learning rate              & 0.001        & Kingma \& Ba (2015) \\
% $\gamma$    & Discount factor                 & 0.95         & Near-future focused \\
% $\varepsilon_0$ & Initial exploration rate    & 0.30         & Standard DQN setting \\
% $\varepsilon_{\min}$ & Minimum exploration   & 0.05         & Maintain exploration \\
% $T_\varepsilon$ & Epsilon decay horizon (steps) & 10,000    & 10 min at 10 tasks/s \\
% $B$         & Mini-batch size                 & 64           & GPU memory trade-off \\
% $|\mathcal{B}|_{\max}$ & Replay buffer size  & 50,000       & Decorrelation coverage \\
% $T_{\text{sync}}$ & Target network sync (steps) & 1,000    & Stable target period \\
% $\delta$    & Huber loss threshold            & 1.0          & Standard DQN value \\
% \midrule
% \multicolumn{4}{l}{\textit{Reward weights}} \\
% $\omega_L$  & Latency weight                  & 0.5          & Primary objective \\
% $\omega_E$  & Energy weight                   & 0.3          & Secondary objective \\
% $\omega_V$  & Violation weight                & 0.2          & Constraint enforcement \\
% $\lambda_{\text{crit}}$ & Safety multiplier  & 10.0         & Safety-critical penalty \\
% $E_{\text{ref}}$ & Energy normalisation ref. (J) & 0.10     & Typical task energy \\
% \midrule
% \multicolumn{4}{l}{\textit{Mobility}} \\
% $R_k$       & Fog coverage radius (m)         & 250          & 5G small cell radius \\
% $T_{\max}$  & T\_exit normalisation bound (s) & 10.0         & Maximum useful lookahead \\
% $s_{\max}$  & Speed normalisation bound (m/s) & 33.3         & 120 km/h maximum \\
% \midrule
% \multicolumn{4}{l}{\textit{Simulation}} \\
% $N$         & Number of vehicles              & 50           & Urban density scenario \\
% $\lambda_{\text{task}}$ & Task arrival rate (Hz) & 10       & 10 fps camera \\
% $T_{\text{sim}}$ & Simulation duration (s)    & 600          & 10-minute evaluation \\
% $T_{\text{warm}}$ & Warm-up period (s)        & 60           & Agent stabilisation \\
% \bottomrule
% \end{tabularx}
% \end{table}

% ================================================================
\section{End-to-End Processing Pipeline}
\label{sec:pipeline}
% ================================================================

The complete processing pipeline for a single DAG task $\tau$ submitted by
vehicle $v_i$ is described by Algorithm~\ref{alg:pipeline}. It integrates
all five algorithmic components into a unified sequential procedure.

\begin{algorithm}[H]
\caption{PCNME End-to-End Task Processing}
\label{alg:pipeline}
\begin{algorithmic}[1]
\State \textbf{Input:} DAG task $\tau$, fog state $\{\rho_k, q_k\}$,
  vehicle state $\sigma_i$, pre-trained agent $Q(\cdot;\,\bm{\theta})$
\State \textbf{Output:} completed task with latency $L_\tau$, energy $E_\tau$,
  deadline verdict
\For{each offloadable step $g_j \in G$}
  \State Compute $\mathrm{EC}(g_j)$ via Eq.~\eqref{eq:ec}
  \If{$\mathrm{EC}(g_j) \geq \theta$}
    \State Route $g_j$ to $\mathcal{C}$ directly \Comment{Boulder}
    \State Execute and record $L_j$, $E_j$; \textbf{continue}
  \EndIf
  \If{$Q(t) > Q_{\max}$}
    \State Trigger Aggregator; bundle overflow pebbles to $\mathcal{C}$
  \EndIf
  \State Build state vector $\mathbf{s}_j$ via Eq.~\eqref{eq:state}
  \State $a_j \leftarrow \varepsilon\text{-greedy}(Q(\mathbf{s}_j;\,\bm{\theta}))$
    \Comment{Agent placement decision}
  \State $x_j \leftarrow \mathcal{X}[a_j]$ \Comment{Map action to destination}
  \State Compute $T_{\text{exit}}(v_i, x_j, t)$ via Eq.~\eqref{eq:texit}
  \State Compute $T_{\text{exec}}(g_j, x_j, t)$ via Eq.~\eqref{eq:exec_fog}
  \State $\mathrm{mode} \leftarrow$ select\_mode via Eq.~\eqref{eq:mode}
  \If{$\mathrm{mode} = \textsc{Proactive}$}
    \State Identify $f_{k'}$ via Eq.~\eqref{eq:dest_node}
    \State Pre-spin container at $f_{k'}$ with resource reservation
    \State Redirect $g_j$ to $f_{k'}$; update $x_j \leftarrow f_{k'}$
  \EndIf
  \State Transmit $g_j$ to $x_j$; record $T_{\text{access}}$
  \State Execute $g_j$ on $x_j$; record $T_{\text{exec}}$
  \State $L_j \leftarrow T_{\text{access}} + T_{\text{exec}}$;
    $E_j \leftarrow$ Eq.~\eqref{eq:step_energy}
  \State $r_j \leftarrow R(g_j, a_j)$ via Eq.~\eqref{eq:reward}
  \State Store transition $(s_j, a_j, r_j, s_{j+1})$ in $\mathcal{B}$
  \State Sample $\mathcal{B}' \subset \mathcal{B}$; update $\bm{\theta}$ via
    Eq.~\eqref{eq:huber}
\EndFor
\State $L_\tau \leftarrow \sum_{g_j \in \Pi} L_j$;
  $E_\tau \leftarrow \sum_j E_j$
\State Report: $L_\tau$, $E_\tau$, $L_\tau \leq D_\tau$
\end{algorithmic}
\end{algorithm}

% ================================================================
\section{Evaluation Metrics}
\label{sec:metrics}
% ================================================================

The following metrics are used to evaluate the framework, each derived from
the formal quantities defined above.

\begin{align}
  \text{Task Feasibility Rate} &=
  \frac{1}{|\mathcal{T}|} \sum_{\tau \in \mathcal{T}}
  \mathbf{1}[L_\tau \leq D_\tau]
  \label{eq:feasibility}\\[6pt]
  \text{Average Latency} &= \frac{1}{|\mathcal{T}|}
  \sum_{\tau \in \mathcal{T}} L_\tau
  \label{eq:avg_latency}\\[6pt]
  \text{Average Energy} &= \frac{1}{|\mathcal{T}|}
  \sum_{\tau \in \mathcal{T}} E_\tau
  \label{eq:avg_energy}\\[6pt]
  \text{Handoff Success Rate} &=
  \frac{\text{handoffs with no re-submission}}{\text{total handoff events}}
  \label{eq:handoff}\\[6pt]
  \text{Pareto Hypervolume} &= \lambda\!\left(
  \bigcup_{\mathbf{c} \in \mathcal{P}^*}
  \bigl\{ \bm{y} \;\big|\; \bm{f}(\mathbf{c}) \preceq \bm{y} \preceq \bm{f}^{\text{ref}} \bigr\}
  \right)
  \label{eq:hv}
\end{align}

where $\lambda(\cdot)$ in~\eqref{eq:hv} denotes Lebesgue measure and
$\bm{f}^{\text{ref}}$ is a reference point dominated by all Pareto solutions.
Hypervolume is a quality indicator for the offline Pareto front that captures
both proximity to the true front and diversity of solutions.

% ================================================================
% REFERENCES
% ================================================================
\begin{thebibliography}{16}

\bibitem{jafari2021}
V.~Jafari and M.\,H.~Rezvani,
``Joint optimization of energy consumption and time delay in {IoT}--fog--cloud
  computing environments using {NSGA-II} metaheuristic algorithm,''
\textit{J. Ambient Intell. Humanized Comput.}, vol.~14, pp.~1675--1698, 2023.

\bibitem{nouri2024tof}
N.~Nouri and M.\,H.~Rezvani,
``A multi-objective approach for optimizing {IoT} applications offloading
  in fog--cloud environments with {NSGA-II},''
\textit{J. Supercomput.}, 2024.

\bibitem{li2025microservice}
Y.~Li et al.,
``Energy-efficient resource management in microservices-based fog and edge
  computing: state-of-the-art and future directions,''
\textit{ACM Comput. Surv.}, 2025.

\bibitem{zhu2025dag}
L.~Zhu et al.,
``Dynamic task offloading scheme for edge computing via
  meta-reinforcement learning,''
\textit{J. Manuf. Syst.}, 2025.

\bibitem{mavto2024}
S.~Zaidi et al.,
``Task offloading optimization using {PSO} in fog computing for the
  internet of drones,''
\textit{Drones}, vol.~8, no.~11, p.~696, 2024.

\bibitem{vfc_drl2024}
M.\,P.~Toopchinezhad and M.~Ahmadi,
``Deep reinforcement learning for delay-optimized task offloading in
  vehicular fog computing,''
\textit{arXiv:2410.03472}, 2024.

\bibitem{paper5_mobility}
M.~Zhu, J.~Liu, and H.~Wang,
``A collaborative offloading task framework for {IoT} fog computing,''
in \textit{Proc. Int. Conf. Fog Mobile Edge Comput.}, 2023.

\bibitem{gasmi2024sdn}
K.~Gasmi and T.~Ezzedine,
``Reinforcement learning based task offloading of {IoT} applications
  in fog computing,''
\textit{Cluster Comput.}, 2024.

\bibitem{zhu2025spatiotemporal}
L.~Zhu et al.,
``A vehicular edge computing offloading and task caching solution
  based on spatiotemporal prediction,''
\textit{Future Gener. Comput. Syst.}, vol.~166, p.~107679, 2025.

\bibitem{deb2002nsga2}
K.~Deb, A.~Pratap, S.~Agarwal, and T.~Meyarivan,
``A fast and elitist multiobjective genetic algorithm: {NSGA-II},''
\textit{IEEE Trans. Evol. Comput.}, vol.~6, no.~2, pp.~182--197, Apr. 2002.

\bibitem{mnih2015dqn}
V.~Mnih et al.,
``Human-level control through deep reinforcement learning,''
\textit{Nature}, vol.~518, pp.~529--533, Feb. 2015.

\bibitem{kingma2015adam}
D.\,P.~Kingma and J.~Ba,
``Adam: a method for stochastic optimization,''
in \textit{Proc. Int. Conf. Learn. Representations (ICLR)}, 2015.

\bibitem{sutton2018rl}
R.\,S.~Sutton and A.\,G.~Barto,
\textit{Reinforcement Learning: An Introduction}, 2nd~ed.
MIT Press, Cambridge, MA, 2018.

\bibitem{huber1964}
P.\,J.~Huber,
``Robust estimation of a location parameter,''
\textit{Ann. Math. Stat.}, vol.~35, no.~1, pp.~73--101, 1964.

\bibitem{mao2016np}
Y.~Mao, C.~You, J.~Zhang, K.~Huang, and K.\,B.~Letaief,
``A survey on mobile edge computing: the communication perspective,''
\textit{IEEE Commun. Surv. Tuts.}, vol.~19, no.~4, pp.~2322--2358, 2017.

\bibitem{mashal2024drl}
I.~Mashal et al.,
``Multiobjective offloading optimization in fog computing using deep
  reinforcement learning,''
\textit{J. Comput. Netw. Commun.}, 2024.

\end{thebibliography}

\end{document}